% from ega0-1-fr.tex
\oldpage[0\textsubscript{I}]{71}

  sur $A/\mathfrak J$. Si on pose $\mathfrak J'=\mathfrak JA'$, $A'$ est alors
  un anneau $\mathfrak J'$-adique, et
  $\mathfrak J'/\mathfrak J'^{\,2}$ est de type fini sur
  $A'/\mathfrak J'$. Si, en outre, $A$ est noethérien, il en est de même de
  $A'$.
\end{enumerate}
\end{proposition}

\begin{enumerate}
  \item[(i)] Si $\mathfrak J$ est un idéal de $A$, $\mathfrak J'$ l'idéal de
  $A'$ formé des $\sum_\alpha c_\alpha T^\alpha$ tels que
  $c_\alpha\in\mathfrak J$ pour tout $\alpha$, alors
  $(\mathfrak J')^n\subset(\mathfrak J^n)'$~; si $\mathfrak J$ est un idéal de
  définition de $A$, $\mathfrak J'$ est donc un idéal de définition de $A'$.

  \item[(ii)] Posons $A_i=A/\mathfrak J^{i+1}$, et pour $i\leq j$, soit
  $u_{ij}$ l'homomorphisme canonique
  $A/\mathfrak J^{j+1}\to A/\mathfrak J^{i+1}$~; posons
  $A'_i=A_i[T_1,\ldots,T_r]$, et soit $u'_{ij}$ l'homomorphisme
  $A'_j\to A'_i$ ($i\leq j$) obtenu en appliquant $u_{ij}$ aux coefficients
  des polynômes de $A'_j$. Montrons que le système projectif
  $(A'_i,u'_{ij})$ vérifie les conditions de (7.2.7)~; comme $\mathfrak J'$ est
  le noyau de $A'\to A'_0$, cela démontrera la première assertion de (ii). Or,
  il est clair que les $u'_{ij}$ sont surjectifs~; le noyau $\mathfrak J'_i$
  de $u'_{0i}$ est l'ensemble des polynômes de
  $A_i[T_1,\ldots,T_r]$ dont les coefficients sont dans
  $\mathfrak J/\mathfrak J^{i+1}$~; en particulier, $\mathfrak J'_1$ est
  l'ensemble des polynômes de $A_1[T_1,\ldots,T_r]$ dont les coefficients
  sont dans $\mathfrak J/\mathfrak J^2$. Comme $\mathfrak J/\mathfrak J^2$
  est de type fini sur $A_1=A/\mathfrak J^2$, on voit déjà que
  $\mathfrak J'_1/\mathfrak J_1'^{\,2}$ est un module de type fini sur $A'_1$
  (ou, ce qui revient au même, sur
  $A'_0=A'_1/\mathfrak J'_1$). Montrons que le noyau de $u'_{ij}$ est
  $\mathfrak J_j'^{\,i+1}$. Il est évident que
  $\mathfrak J_j'^{\,i+1}$ est contenu dans ce noyau. D'autre part, soient
  $a_1,\ldots,a_m$ des éléments de $\mathfrak J$ dont les classes
  mod.~$\mathfrak J^2$ engendrent $\mathfrak J/\mathfrak J^2$~; on vérifie
  immédiatement que les classes mod.~$\mathfrak J^{j+1}$ des monômes de degré
  $\leq j$ en les $a_k$ ($1\leq k\leq m$) engendrent
  $\mathfrak J/\mathfrak J^{j+1}$, et les classes des monômes de degré $>i$
  et $\leq j$ engendrent donc
  $\mathfrak J^{i+1}/\mathfrak J^{j+1}$~; un monôme en les $T_k$ ayant un tel
  élément comme coefficient est donc produit de $i+1$ éléments de
  $\mathfrak J'_j$, ce qui établit notre assertion. Enfin, si $A$ est
  noethérien, il en est de même de
  $A'/\mathfrak J'=(A/\mathfrak J)[T_1,\ldots,T_r]$, donc $A'$ est noethérien
  (7.2.8).
\end{enumerate}

\begin{proposition}[7.5.5]
\label{0.7.5.5-fr}
Soient $A$ un anneau noethérien $\mathfrak J$-adique, $B$ un anneau topologique
admissible, $\varphi:A\to B$ un homomorphisme continu, faisant de $B$ une
$A$-algèbre. Les conditions suivantes sont équivalentes~:
\begin{enumerate}
  \item[a)] $B$ est noethérien et $\mathfrak JB$-adique, et
  $B/\mathfrak JB$ est une algèbre de type fini sur $A/\mathfrak J$.

  \item[b)] $B$ est topologiquement $A$-isomorphe à $\varprojlim B_n$, où
  $B_n=B_m/\mathfrak J^{n+1}B_m$ pour $m\geq n$, et $B_1$ est une algèbre de
  type fini sur $A_1=A/\mathfrak J^2$.

  \item[c)] $B$ est topologiquement $A$-isomorphe à un quotient d'une algèbre
  de la forme\linebreak $A\{T_1,\ldots,T_r\}$ par un idéal (nécessairement fermé en vertu
  de (7.3.6) et (7.5.4 (ii))).
\end{enumerate}
\end{proposition}

Comme $A$ est noethérien, il en est de même de
$A'=A\{T_1,\ldots,T_r\}$ (7.5.4), donc c) implique que $B$ est noethérien~;
comme $\mathfrak J'=\mathfrak JA'$ est un voisinage ouvert de 0 dans $A'$ tel
que les $\mathfrak J'^{n}$ forment un système fondamental de voisinages de 0,
les images $\mathfrak J^nB$ des $\mathfrak J'^{n}$ forment un système
fondamental de voisinages de 0 dans $B$, et comme on sait que $B$ est séparé et
complet, $B$ est un anneau $\mathfrak JB$-adique. Enfin, $B/\mathfrak JB$ est
une algèbre (sur $A/\mathfrak J$) quotient de
$A'/\mathfrak JA'=(A/\mathfrak J)[T_1,\ldots,T_r]$, donc est de type fini, ce
qui achève de prouver que c) entraîne a).

Si $B$ est $\mathfrak JB$-adique et noethérien, $B$ est isomorphe à
$\varprojlim B_n$, où $B_n=B/\mathfrak J^{n+1}B$ (7.2.11), et
$\mathfrak JB/\mathfrak J^2B$ est un module de type fini sur
$B/\mathfrak JB$. Soit $(a_j)_{1\leq j\leq s}$ un système de générateurs du
$(B/\mathfrak JB)$-module $\mathfrak JB/\mathfrak J^2B$, et
$(c_i)_{1\leq i\leq r}$ un système d'éléments de $B/\mathfrak J^2B$ dont les
classes

